% -*-coding: utf-8 -*-
%%%%%%%%%%%%%%%%%%%%%%%%%%%%%%%%%%%%%%%%%%%%%%%%%%%%%%%%%
\appendix
\chapter{协同SFADA风险上界证明}[协同SFADA风险上界证明]%
本附录给出第 3 章定理 3.4的完整证明。引理 3.1 与引理 3.2 给出的 UDA 与 SFDA 风险上界作为前置结论，其证明可见标准域适应理论文献~\cite{mathelin2022discrepancy}。
%%%%%%%%%%%%%%%%%%%%%%%%%%%%%%%%%%%%%%%%%%%%%%%%%%%%%%%%%
% \section{定理 3.4 的证明}[Proof of Theorem 3.4]
\label{app:theorem_3_4_proof}

\begin{proof}
由 $L$ 关于各变量的 $\mu$-Lipschitz 连续性与三角不等式，对所有 $x'\in \mathcal T$，$x\in \mathcal T_q$，目标模型风险可分解为
\begin{align}
\mathcal L_P(h,f_P) \ &\le\ \underbrace{\mathbb E_{x\sim P}L\big(f_P(x),f'_P(x)\big)}_{(A)}  + \underbrace{\mathbb E_{x\sim P}\!\big[L(f'_P(x),g(x))+L(g(x),h(x))\big]}_{(B)}.
\end{align}

\paragraph{控制项 (B)}
令 $g^\star\in\arg\min_g\max_{x\in\mathcal T}\big[L(f'_P(x),g(x))+L(g(x),h(x))\big]$，并记 $F_g(x):=L(f'_P(x),g(x))+L(g(x),h(x))$，则对任意 $x\in \mathcal T$ 有
\begin{equation}
F_g(x)\le \max_{u\in\mathcal T} F_g(u).
\end{equation}
按 $\eta_{\mathcal G}$ 的定义：
\begin{equation*}
\max_{u\in\mathcal T}F_{g^\star}(u)=\eta_{\mathcal G}.
\end{equation*}
对 $x\sim P$ 取期望并应用 Hoeffding / McDiarmid 不等式与对 $h$、$g$ 的联合 Rademacher 复杂度一致界，以至少 $1-\delta/3$ 的概率：
\begin{align}
(B)(g)=\mathbb E_{x\sim P}F_g(x) &\le \frac{1}{n_t}\sum_{x'\in \mathcal T}F_g(x') + 2\mu \big(\mathfrak R_{n_t}(\mathcal H)+ \mathfrak R_{n_t}(\mathcal G \circ\phi)\big)+\sqrt{\tfrac{M^2\log(3/\delta)}{2n_t}}.
\end{align}
将 $g^\star$ 代入，并利用 $\frac{1}{n}\sum_{x'\in \mathcal T}F_g(x')\le \max_{u\in\mathcal T}F_{g^\star}(u)=\eta_{\mathcal G}$，得
\begin{align}
(B) \ &\le\ \eta_{\mathcal G} \ +\ 2\mu\big(\mathfrak R_{n_t}(\mathcal H)+\mathfrak R_{n_t}(\mathcal G\circ\phi)\big) +\sqrt{\tfrac{M^2\log(3/\delta)}{2n_t}}.
\end{align}

\paragraph{控制项 (A)}
由经验均值的集中不等式，以至少 $1-\delta/3$ 的概率：
\begin{align}
(A) \ &\le\ \frac{1}{n_t}\sum_{x'\in\mathcal T}L\big(f_P(x'),f'_P(x')\big) \ +\ \sqrt{\tfrac{M^2\log(3/\delta)}{2n_t}}.
\end{align}
对 $x\in \mathcal T_q$，由引理3.3 并对 $x'$ 取平均：
\begin{align}
   \frac{1}{n_t}\sum_{x'\in\mathcal T} L\big(f_P(x'),f'_P(x')\big) &\le \frac{1}{|\mathcal T_q|}\sum_{x\in\mathcal T_q} L\big(f_P(x),f'_P(x)\big)\notag \\
   &\quad+\frac{\mu}{n_t}\sum_{x'\in\mathcal T}\Big(k d_X(x',x)+ k_V d_Z\big(\phi(x'),\phi(x)\big)\Big).
\end{align}
由于上式对任意 $x\in\mathcal T_q$ 成立，特别地可取 $x\in\mathcal T_q$ 使 $x$ 尽量接近 $x'$，从而
\begin{align}
   \frac{1}{n_t}\sum_{x'\in\mathcal T} L\big(f_P(x'),f'_P(x')\big) &\le \frac{1}{|\mathcal T_q|}\sum_{x\in\mathcal T_q} L\big(f_P(x),f'_P(x)\big)\notag \\
   &\quad+\frac{\mu}{n_t}\sum_{x'\in\mathcal T}\Big(k d_X(x',\mathcal T_q)+ k_V d_Z\big(\phi(x'),\mathcal T_q\big)\Big).
\end{align}

合并 (A) 与 (B) 的上界，并对三个 $\delta/3$ 事件取并集，以至少 $1-\delta$ 的概率：
\begin{align}
    \mathcal L_{P}(h,f) &\le\mathcal L_{\hat Q}(f_P,f'_P)+\frac{\mu}{n_t}\sum_{x'\in\mathcal T}\!\Big(k d_X(x',\mathcal T_q)+k_V d_Z(\phi(x'),\mathcal T_q)\Big) \notag\\
    &\quad+\eta_{\mathcal G}+2\mu\big(\mathfrak R_{n_t}(\mathcal H)+\mathfrak R_{n_t}(\mathcal G\!\circ\!\phi)\big) +2\sqrt{\tfrac{M^2\log(3/\delta)}{2n_t}}.
\end{align}
即式~\eqref{eq:collab_bound}，定理3.4 证毕。
\end{proof}

